\documentclass[a4paper,11pt]{ltjsarticle}
\usepackage[haranoaji]{luatexja-preset}
\usepackage{amsmath,amssymb}
\usepackage[top=12mm,bottom=10mm,left=14mm,right=14mm]{geometry}
\usepackage{array}
\usepackage{tikz}
\usetikzlibrary{calc}
\pagestyle{empty}
\setlength{\parindent}{0pt}
\newlength{\qcol}\setlength{\qcol}{0.47\textwidth}
\newlength{\qgap}
\newlength{\anscolw}\setlength{\anscolw}{0.30\textwidth}
\newlength{\ansh}
\newcommand{\Q}[2]{\parbox[t]{\qcol}{\raggedright (#1)\hspace{0.6em}$\displaystyle #2$}}
\newcommand{\QR}[4]{\Q{#1}{#2}\hfill\Q{#3}{#4}\par\vspace{\qgap}}
\newcommand{\anscell}[1]{\rule[-\ansdrop]{0pt}{\ansh}(#1)}
\newlength{\ansdrop}

\begin{document}
{\large\bfseries 数学テスト　20}\hfill 名前(\underline{\hspace{32mm}})\par\vspace{3mm}
■（1）〜（12）の計算をしなさい。（13）、（14）は連立方程式を解きなさい。\par\vspace{3mm}
\setlength{\qgap}{11mm}
\QR{1}{9x^2y\div(-3xy)}{2}{-5(2a-b)}
\QR{3}{(7a^2-4a-1)-(2a^2-3a)}{4}{(5x-4y)+(-5x+4y)}
\QR{5}{(12x-6y+24)\div6}{6}{4a^2\div(-5b)\times10b^2}
\QR{7}{-48xy^3\div4xy^2}{8}{-2xy\times6yz}
\QR{9}{2x^2y\times3xy^2\div\left(-\dfrac{1}{2}x^2y\right)}{10}{\dfrac{x-3y}{8}+\dfrac{2x+y}{4}}
\QR{11}{-8(x+2y)+7(2x+3y)}{12}{\dfrac{1}{4}(8a+12b)-\dfrac{2}{3}(18a-6b)}
\QR{13}{\left\{\begin{array}{l} \dfrac{2x-3}{4}+\dfrac{y+9}{8}=1 \\ -0.3x-0.1y=0.2 \end{array}\right.}{14}{\left\{\begin{array}{l} 3y=2x+5 \\ x+6y=5 \end{array}\right.}
\vfill
\setlength{\ansh}{12mm}\setlength{\ansdrop}{8mm}%
\begin{center}\begin{tabular}{|p{\anscolw}|p{\anscolw}|p{\anscolw}|}\hline
\anscell{1} & \anscell{2} & \anscell{3} \\\hline
\anscell{4} & \anscell{5} & \anscell{6} \\\hline
\anscell{7} & \anscell{8} & \anscell{9} \\\hline
\anscell{10} & \anscell{11} & \anscell{12} \\\hline
\anscell{13} & \anscell{14} &  \\\hline
\end{tabular}\end{center}

\end{document}
